Una vareta metàl·lica es desplaça a una velocitat constant de $3,00\ \mathrm{m\,s^{-1}}$ sobre un conductor en forma de U dins un camp magnètic uniforme $B = 0,40\ \mathrm{T}$, perpendicular al pla i en el sentit indicat a les figures. En l'instant inicial, la vareta es troba a una distància de 3,00~cm de la part esquerra del conductor.

\begin{apartat}{1,25}
Calculeu el flux magnètic a través del circuit tancat per la vareta en l'instant inicial indicat a la figura~1. En el moment en què la vareta es posa en moviment, apareix un corrent induït en el sentit indicat a la figura~2. Justifiqueu quin és el sentit de moviment de la vareta, a l'esquerra o a la dreta, i dibuixeu el vector de velocitat sobre la figura de la vareta en moviment. Determineu la FEM induïda mentre la vareta es mou a $3,00\ \mathrm{m\,s^{-1}}$.
\end{apartat}

\figura[0.75]{fig1.png}

\begin{apartat}{1,25}
Considereu la situació il·lustrada a la figura~2, en què la intensitat induïda és de 6,00~mA. Justifiqueu si en aquestes condicions apareixerà alguna força magnètica sobre la vareta. En cas afirmatiu, calculeu el mòdul de la força magnètica sobre la vareta i dibuixeu-la.
\end{apartat}
